\section{模型假设}

（1）药材内部的温度与水分只沿径向变化，忽略轴向与周向的不均匀。药材长度 $25$~cm 与直径
$4$~cm 之比为 $6.25$，两端面面积仅占总换热面积的 $7.4\%$；按扩散时间 $L^2/D$ 估计，
轴向输运所需时间约为径向的 $39$ 倍，在本题关心的时间范围内轴向梯度不足以改变径向结果。

（2）药材视为均质连续介质，水分以有效扩散的方式迁移，不区分液态水与水蒸气的输运路径。
题目直接给出有效水分扩散系数的表达式，即已把毛细流动与气相扩散归并到同一个系数中，
这是干燥文献中处理多孔物料的常用简化\upcite{Srikiatden2007}。

（3）能量方程不含水分蒸发的潜热汇。题目未给出汽化潜热与相变发生位置，无法确定该汇项的
大小与分布。后文采用近似汽化潜热作结构性诊断，以评估该省略对时长的影响方向与量级。

（4）烘房空气温度与含湿量在 $4$~h 后采用稳态外推值 $50.0000~^\circ$C 与 $0.0500$~kg/kg，不作线性外推。附件 1 的
温度在约 $2.95$~h 后已进入平台，继续线性外推会得到无物理依据的持续升温。

（5）药材收缩只发生在径向，长度不变，且收缩不可逆。题目仅给出半径随时间的实测数据，
未给出长度变化；失水造成的骨架塌缩在恒温干燥中不会自行回复，故半径取为单调不增。

（6）表面对流换热系数与对流传质系数在四个子问题中取同一组给定值，不随温度、含水率或
表面收缩而变化。题目在给出这两个系数时未附带修正关系，按题面条件统一使用。
